
\section*{初版にむけたあとがき}

2020年の2月から始まったコロナ感染症によるさまざまな社会変化は，2021年度からの授業スタイルも大きく変えてしまいました。
大学はオンラインで授業する方法を考え，教職員にその対応を支持しなければなりませんでしたし，我々教員の方も一人ひとりが授業の質を落とさず実践できるかを考えさせられたのです。

このテキストはそのような混乱の中で生まれました。
それまでの授業は，必要な参考文献を随時参照しながら，自作のスライドを使って講義をするスタイルにしていましたが，そのやり方では限界があるのではないかと思われました。
学生がキャンパスに入れないのですから，何冊かの本をあちこち参照するスタイルが難しい，というのがまず1つ。
学生の顔を見て授業の進捗を調節するなど，授業フィードバックを通じて補足したりするスタイルが難しいというのもあります。オンライン対応が始まった当初は受講生の通信環境を考え，リアルタイム動画通信など，通信データサイズが大きくなるのは避けるようにという方針があったからです。

実はコロナの問題に関係なく，授業の計画を詳細なものにしようという取り組みは始めていました。\textcite{Ashida2019}によってコマシラバス，詳細な授業設計の重要性に気付かされてから，授業全体を見通した設計をし直そうと考えたからです。それの準備を始めていた頃，コロナ問題がやってきました。授業計画は新年度前に確立できたものの，このテキストのような資料は準備ができていませんでしたし，前期後期15回ずつの授業実施計画も叶いませんでしたので，即活用できたということにはなりませんでしたが。

またそれまでのスライド教材を見直すと，箇条書きになっている項と項の間に話していたことも，不均一で不完全であるように思われました。それまでもまじめに授業資料を作り込んできたつもりでしたが，よく読むと「これでは伝わらない」「これでは行間に暗黙の仮定や文脈に依存した情報が多すぎる」など，数多くの気になる点が出てきてしまったのです。そこで，前期には間に合いませんでしたが，後期には授業で話す内容をすべて文字に書き起こすという教材作りが始まりました。
このテキストは，2020年後期から2021年前期を通じて用意された書き起こし教材を元に，修正や調整を加えて作成されたものです。またこのテキストには行番号が付されてありますが，これは授業が「どこを参照しているか」が明確にわかるような仕掛けであり，\textcite{Ashida2019}に含まれているアイデアの1つです。

このテキストはバージョン管理されています。バージョンはPDF版でしたら各ページの右上に，KDP版でしたら奥付けに書いてあります。バージョン番号は，メジャー，マイナー，パッチの3つの要素から構成されています。
今はメジャーバージョン2です。シラバスの骨組みが変わるようなことがあれば，これが1つ増えることになります。チャプターレベルの変更があればマイナーの数字が上がります。
誤字脱字など，細かな変更はパッチの数字が変わります。なんせ自作のテキストですから，誤字脱字や誤変換，数字や参照のミスが多く，その度に受講生に指摘してもらってはアップデートする，ということを繰り返してここまでやってきました。

最新版のシラバスやテキストは，サイト\footnote{\url{https://kosugitti.github.io/psychometrics_syllabus/}}で公開されています。
GitHubを使って管理し，図版のほとんどはRを使って作っています。版組にはもちろんLua\LaTeX を使っています。一番最近のアップデートは，引用文献をBib\LaTeX を使ったものにしたことでしょうか。
バージョンがこまめに上がるものですし，リンク機能などPDFならではの利点もあるので，できれば電子デバイスで電子ファイルのまま利用してもらえればと思っています。
しかし受講生の多くはプリントアウトして使っているようですから，この度実費で製本したものが入手できるよう，Kindle Direct Publishingを利用してペーパーバック版も用意しました。
専修大学人間科学部心理学専攻の学生向けに用意したテキストではありますが，どなたでもお読みいただけるようになっています。内容についてはできるだけの注意を払って，誤解や間違いがないよう気をつけてはおりますが，おかしなところがありましたら著者までご連絡ください。

また原稿そのものはCreative Commons BY-SA(CC BY-SA)ライセンスVersion 4.0に基づいて提供しています。加筆修正や再頒布など，ライセンスに従ってご活用いただければと思います。


不十分で不完全なところもあるかと思いますが，この取り組みが少しでも，受講生の皆さんのお役に立てば幸いです。

\begin{flushright}
	\date{\today}\\
	小杉考司
\end{flushright}

